\documentclass[11pt,a4paper]{article}
\usepackage[utf8]{inputenc}
\usepackage[spanish]{babel}
\usepackage{amsmath,amsfonts,amssymb}
\usepackage{geometry}
\geometry{top=2.5cm,bottom=2.5cm,left=2.5cm,right=2.5cm}
\usepackage{hyperref}
\usepackage{booktabs}
\usepackage{xcolor}
\usepackage{tcolorbox}
\usepackage{enumitem}
\usepackage{tikz}
\usetikzlibrary{shapes.geometric, arrows, positioning, calc}

% Definición de colores institucionales
\definecolor{externadoDarkBlue}{RGB}{10, 34, 64}
\definecolor{externadoAccent}{RGB}{0, 114, 206}
\definecolor{lightGray}{RGB}{245, 247, 250}
\definecolor{alertRed}{RGB}{180, 40, 40}

\hypersetup{
    colorlinks=true,
    linkcolor=externadoAccent,
    urlcolor=externadoAccent,
    titlecolor=externadoDarkBlue
}

\tcbset{
    colback=lightGray,
    colframe=externadoDarkBlue,
    fonttitle=\bfseries\color{white},
    arc=3mm
}

\title{
    \vspace{-1cm}
    \textbf{\Large UNIVERSIDAD EXTERNADO DE COLOMBIA}\\
    \textbf{\large PREGRADO EN CIENCIA DE DATOS --- EMPRENDIMIENTO}\\
    \vspace{0.5cm}
    \Large \textbf{TALLER 3: DIAGNÓSTICO ESTRATÉGICO, REENCUADRE Y FEEDBACK}\\
    \large \textbf{Proyecto: StatCredit AI --- Motor SaaS B2B de Scoring Crediticio Alternativo \& XAI}
}
\author{\textbf{Grupo 5:} Santiago Sandoval | Sebastián Ramos | Julián Duarte | Tomás Rincón | Julián Jiménez}
\date{Semestre 2026-II}

\begin{document}

\maketitle

\hrule
\vspace{0.5cm}

\section{Introducción y Contexto Estratégico}
El presente documento consolida el diagnóstico estratégico refinado de \textbf{StatCredit AI}, incorporando las observaciones de evaluación cualitativa sobre la estructuración del problema, la economía informal, la teoría del credit scoring y la distinción de nuestro modelo de negocio SaaS B2B.

StatCredit AI aborda la asimetría de información que impide a las entidades financieras tradicionales originar crédito de forma segura hacia más de 13 millones de trabajadores independientes y micronegociantes solventes en Colombia.

---

\section{Reencuadre Estratégico del Problema (Bianchi \& Verganti, 2021)}
Siguiendo la metodología de \textit{Problem Framing \& Reframing} (\textit{Bianchi \& Verganti, 2021}), la innovación disruptiva surge al modificar el marco de referencia interpretativo de una problemática:

\begin{tcolorbox}[title=El Reencuadre de StatCredit AI]
\begin{itemize}[leftmargin=*]
    \item \textbf{Problema Visible Tradicional:} \textit{«Los trabajadores independientes y micronegocios son clientes de alto riesgo moroso porque carecen de historial crediticio formal en centrales de riesgo (DataCrédito/TransUnion).»}
    \item \textbf{Reencuadre StatCredit AI:} \textit{«Los trabajadores independientes poseen una intensa actividad financiera y transaccional diaria en billeteras digitales y comercio no bancario, pero los modelos lineales tradicionales carecen de la capacidad metodológica e infraestructura para interpretar dicha información como un activo de repago.»}
    \item \textbf{Enunciado Insignia:} \textit{“No buscamos cambiar quién puede acceder al crédito; buscamos cambiar la forma en que se mide quién puede pagarlo.”}
\end{itemize}
\end{tcolorbox}

\subsection{Los 4 Movimientos de Diseño de Problemas}
\begin{enumerate}
    \item \textbf{De Déficit a Potencial:} El independiente no carece de capacidad de pago; posee una dinámica de caja no canalizada por los formatos bancarios tradicionales.
    \item \textbf{De Problema a Oportunidad:} La banca no rechaza por malicia, sino por falta de herramientas actuariales; existe un mercado desatendido de más de \$45M USD en potencial de originación.
    \item \textbf{De Intuición a Evidencia:} Reemplazamos la discrecionalidad del analista por algoritmos de Inferencia Causal e IA Explicable (XAI).
    \item \textbf{De Solución Aislada a Arquitectura de Plataforma B2B:} Nos posicionamos como el motor inteligente (SaaS/API) que conecta la economía informal con las entidades financieras.
\end{enumerate}

---

\section{Diagrama Causal de Ishikawa (Espina de Pescado) Reestructurado}
Atendiendo las observaciones del cuerpo docente, la **cabeza del problema raíz** radica en la falta de información y datos estructurados en las entidades bancarias para la evaluación del riesgo, mientras que la **exclusión crediticia** y el **gota a gota** son consecuencias de segundo orden.

\begin{figure}[h!]
\centering
\small
\begin{tikzpicture}[node distance=1.5cm, >=stealth]
    % Estilo de cajas
    \tikzstyle{head} = [rectangle, draw=externadoDarkBlue, fill=externadoDarkBlue, text=white, font=\bfseries, text width=4.5cm, align=center, minimum height=1.2cm, rounded corners=2mm]
    \tikzstyle{bone} = [rectangle, draw=externadoAccent, fill=lightGray, font=\bfseries, text width=3.2cm, align=center, rounded corners=1mm]
    \tikzstyle{sub} = [font=\tiny, text width=2.8cm, align=left]

    % Espina central
    \draw[ultra thick, externadoDarkBlue, ->] (0,0) -- (10,0);
    
    % Cabeza (Problema Raíz)
    \node[head] at (11.5,0) {Falta de información financiera estructurada en los bancos para evaluar el riesgo en independientes};

    % Causa Top 1: Datos
    \node[bone] (c1) at (2.5, 2.5) {1. Fragmentación de Datos};
    \draw[thick, externadoAccent, ->] (c1) -- (3.5, 0);
    \node[sub] at (1.8, 1.5) {• Billeteras digitales aisladas\\• Falta de score unificado\\• Datos no estructurados};

    % Causa Top 2: Algoritmos
    \node[bone] (c2) at (7, 2.5) {2. Inadecuación Algorítmica};
    \draw[thick, externadoAccent, ->] (c2) -- (8, 0);
    \node[sub] at (6.3, 1.5) {• Modelos lineales rígidos\\• Penalización de ingresos estacionales\\• Ausencia de inferencia causal};

    % Causa Bottom 1: Economía Informal
    \node[bone] (c3) at (2.5, -2.5) {3. Economía Informal};
    \draw[thick, externadoAccent, ->] (c3) -- (3.5, 0);
    \node[sub] at (1.8, -1.5) {• >50\% ocupación informal\\• Transacciones en Nequi/Efectivo\\• Sin soportes tributarios tradicionales};

    % Causa Bottom 2: Regulación y Cultura
    \node[bone] (c4) at (7, -2.5) {4. Regulación y SARC};
    \draw[thick, externadoAccent, ->] (c4) -- (8, 0);
    \node[sub] at (6.3, -1.5) {• Modelos caja negra rechazados\\• Exigencia de explicabilidad SARC\\• Sesgo pro-asalariado formal};

\end{tikzpicture}
\caption{Diagrama de Ishikawa: Causas raíz del problema de información en el riesgo crediticio.}
\end{figure}

\begin{tcolorbox}[colback=white,colframe=alertRed,title=Consecuencias del Problema Raíz (Efectos de Segundo Orden)]
La falta de información estructurada genera tres consecuencias severas en el ecosistema:
\begin{enumerate}
    \item \textbf{Exclusión y Sobrecosto Financiero B2C:} Los independientes son empujados hacia modalidades informales violentas (\textit{gota a gota}) con tasas $>500\%$ EA.
    \item \textbf{Canibalización y Pérdida de Margen B2B:} Los bancos compiten ferozmente por los mismos 6 millones de asalariados formales.
    \item \textbf{Muerte Prematura de Micronegocios:} El 92\% de las unidades productivas carecen de liquidez para insumos por falta de crédito productivo.
\end{enumerate}
\end{tcolorbox}

---

\section{Fundamentación Teórica y Ejemplo Comparativo del Credit Score}

\subsection{¿Qué es un Score Crediticio?}
Un \textit{Credit Score} es una métrica probabilística estandarizada (rango 150 a 950 puntos) que cuantifica la probabilidad de que un deudor incurra en mora severa ($\ge 90$ días) en un horizonte de 12 meses.

\subsection{Ejemplo Comparativo de Cálculo: Tradicional vs. StatCredit AI}
Para ilustrar la diferencia metodológica y el valor agregado de StatCredit AI, consideremos el caso de un comerciante independiente que factura \$8.000.000 COP mensuales vía Nequi/Daviplata pero carece de tarjeta de crédito bancaria:

\begin{enumerate}[leftmargin=*]
    \item \textbf{Cálculo Tradicional (Regresión Logística en Buró):}
    $$\text{Score}_{\text{Tradicional}} = \beta_0 + \beta_1(\text{HistorialBancario}) + \beta_2(\text{TarjetasCredito}) - \beta_3(\text{ReportesMora})$$
    Dado que $\text{HistorialBancario} = 0$ y $\text{TarjetasCredito} = 0$, el modelo tradicional calcula un puntaje de **320 puntos** ($\implies$ **RECHAZO AUTOMÁTICO DE SOLICITUD**).
    \item \textbf{Cálculo StatCredit AI (Machine Learning No Lineal + XAI):}
    $$\text{Score}_{\text{StatCredit}} = f\big(\text{FlujoTransaccionalNequi}, \text{FrecuenciaDepósitos}, \text{EstabilidadCaja}, \text{FacturaciónDIAN}\big)$$
    El modelo evalúa 12 meses de entradas transaccionales continuas, identifica un patrón de liquidez solvente y asigna **740 puntos** ($\implies$ **APROBACIÓN CON CUPO PROPORCIONAL DE \$5.000.000 COP**).
\end{enumerate}

\subsection{¿Por qué los Bancos Tradicionales No Lo Han Logrado?}
\begin{itemize}[leftmargin=*]
    \item \textbf{Incompatibilidad de Infraestructura de IT:} Los sistemas legacy bancarios están diseñados para procesamiento batch nocturno de archivos planos, no para ingesta vía API en tiempo real de billeteras digitales.
    \item \textbf{Aversión Regulatoria al SARC:} La Superintendencia Financiera (Circular 026) exige auditar por qué se aprueba o rechaza cada crédito. Los bancos temen usar modelos complejos por considerarlos "cajas negras" impredecibles.
    \item \textbf{Sesgo Institucional de Descuento por Nómina:} La banca tradicional prefiere otorgar crédito de consumo a asalariados con deducción automática de nómina, ignorando la oportunidad del sector independiente.
\end{itemize}

\subsection{Ausencia de Evidencia vs. Evidencia de Ausencia}
Los bancos confunden la *falta de reporte en DataCrédito* con *evidencia de insolvencia*. En Colombia, más del 50\% de la fuerza laboral opera en la **economía informal y semi-informal**, moviendo flujos transaccionales reales. StatCredit AI transforma esa ausencia de evidencia en prueba matemática objetiva de capacidad de repago.

---

\section{Caracterización de la Población Objetivo y Modelo B2B SaaS}

\subsection{Población Objetivo End-User (B2B2C)}
Para evitar ambigüedades sobre a quién se dirige la solución, definimos explícitamente nuestra **población objetivo beneficiaria**:
\begin{quote}
\textbf{Población Objetivo:} \textit{Trabajadores independientes solventes de la economía informal y semi-informal urbana en Colombia} (comerciantes, prestadores de servicios autónomos y micronegociantes no bancarizados con facturación o ingresos transaccionales digitales recurrentes entre 2 y 10 SMMLV).
\end{quote}

\subsection{Segmentación Comercial (B2B vs. B2C)}
\begin{table}[h!]
\centering
\small
\begin{tabular}{p{4.5cm} p{5.5cm} p{5.5cm}}
\toprule
\textbf{Dimensión} & \textbf{Cliente Comercial Directo (B2B)} & \textbf{Población Objetivo Beneficiaria (B2B2C)} \\
\midrule
\textbf{Sujeto} & Entidades Financieras, Neobancos, Fintechs de originación, Cooperativas. & Trabajadores independientes, profesionales autónomos y micronegocios urbanos. \\
\midrule
\textbf{Relación Comercial} & Suscriptor pagador de la plataforma SaaS y la API de scoring. & Solicitante de crédito ante la entidad financiera. \\
\midrule
\textbf{Propuesta de Valor} & Reducción de tasa de rechazo, ampliación de colocación y cumplimiento SARC. & Acceso a crédito productivo justo con tasas de interés reguladas. \\
\bottomrule
\end{tabular}
\caption{Matriz de diferenciación B2B vs B2C para StatCredit AI.}
\end{table}

\subsection{Modelo Software as a Service (SaaS)}
StatCredit AI opera como un motor tecnológico **Score-as-a-Service**:
\begin{enumerate}
    \item \textbf{Integración vía REST API / gRPC:} Conexión directa al motor de originación de la entidad B2B.
    \item \textbf{Ingesta de Datos en Tiempo Real:} Recepción en tiempo real de extractos digitales y datos transaccionales con autorización previa.
    \item \textbf{Inferencia y Explicabilidad:} Entrega en <2 segundos del Score (150-950), probabilidad de mora y reporte de explicabilidad local SHAP para auditoría SARC.
    \item \textbf{Estructura de Ingresos:} Suscripción mensual base (mantenimiento MLOps) + Tarifa por consulta (\textit{pay-per-query}).
\end{enumerate}

---

\section{Análisis de la Competencia Actual}

\begin{table}[h!]
\centering
\tiny
\begin{tabular}{p{2.2cm} p{2.2cm} p{2.8cm} p{2.5cm} p{3.2cm}}
\toprule
\textbf{Competidor} & \textbf{Modelo} & \textbf{Producto Principal} & \textbf{Limitación Principal} & \textbf{Diferenciador StatCredit AI} \\
\midrule
\textbf{DataCrédito Experian} & Buró Crediticio Tradicional & Score DataCrédito (150-950) & Punto ciego total con independientes no bancarizados. & Ingesta de datos transaccionales alternativos en tiempo real. \\
\midrule
\textbf{TransUnion} & Buró Crediticio Tradicional & Score TransUnion \& Fraud Check & Optimizado para historial crediticio formal estructurado. & Inferencia Causal para independientes con ingresos estacionales. \\
\midrule
\textbf{Belvo} & Open Finance Aggregator (B2B) & Belvo Data API (Agregación) & Tubería de datos pura; no genera scoring ni asume lógica de riesgo. & Motor inteligente decisorio que entrega un Score predictivo auditable. \\
\midrule
\textbf{Sempli / Addi} & Fintech Originador Directo (Lender) & Crédito Digital Directo & Asumen riesgo de balance y capital de mora. & Proveedor B2B SaaS tecnológico sin riesgo de balance (margen >75\%). \\
\midrule
\textbf{Bancos Tradicionales} & Entidades Financieras & Créditos de Nómina y Consumo & Rechazo masivo de independientes por normas rígidas. & \textbf{Son clientes B2B objetivo}, no competidores directos de la API. \\
\bottomrule
\end{tabular}
\caption{Matriz cualitativa de diferenciación competitiva.}
\end{table}

---

\section{Metodología para la Estimación de Mercado (TAM, SAM, SOM)}

Para el próximo entregable cuantitativo, la metodología se basa en un enfoque \textit{Bottom-Up}:
\begin{itemize}[leftmargin=*]
    \item \textbf{TAM (Total Addressable Market):} Mercado total de soluciones B2B de Credit Scoring y Credit Analytics en América Latina ($\approx \$220\text{M USD}$ anuales).
    \item \textbf{SAM (Serviceable Addressable Market):} Total de solicitudes de crédito originadas por la población independiente y micronegociantes en Colombia ($\approx 13\text{M de personas}$, representando un mercado potencial de scoring de $\$45\text{M USD}$ anuales).
    \item \textbf{SOM (Serviceable Obtainable Market):} Cuota capturable en 3 años mediante alianzas estratégicas con las 25 principales Fintechs de crédito y Cooperativas en Colombia ($\approx \$3.5\text{M USD}$ en ingresos anuales recurrentes ARR).
\end{itemize}

\end{document}
